\Needspace{13\baselineskip}
\section{R1: Search control and shared metavariables}\label{sec:search}

\Q{R1.Q1}{Search without duplicating all dependent work}{Can shared-metavariable goals be organized more effectively than copying entire alternatives or solving every goal sequentially? Can independence be detected and reused?}

\answer Yes: use a persistent whole-state frontier, but represent constraints inside each state as a dependency hypergraph. Independence should be a certified special case of this representation, not the default assumption. Aesop supplies an established best-first, rule-indexed reference design; it does not make arbitrary shared obligations independent.\cite{aesop}

The vertices are outstanding term and universe metavariables. A factor contains an obligation together with every variable on which its meaning can depend. The support must include the target's type, types and values of relevant local declarations, delayed assignments, deferred instance problems, and residual contract observations. A contract $f(0)+g(0)=1$, for example, couples the bodies of $f$ and $g$ even when their Lean types do not share a metavariable. A common immutable local parameter does not itself couple two searches; a mutable variable in its type can.

After an assignment, instantiate the affected factors and update their supports. Components can split when a shared variable becomes fixed. A plain union--find structure is therefore insufficient on its own: it efficiently merges but does not represent this splitting. Recompute small affected components, or use a persistent adjacency representation with a work queue. Start with conservative overapproximations. False dependency edges cost performance; missed dependency edges can invalidate independent pruning.

\begin{proposition}[Factorization criterion]\label{prop:factor}
Suppose the remaining completion variables partition into $V_1,\ldots,V_k$. Each variable's domain, including typing and grammar restrictions, depends only on variables in its own block and fixed parameters. Suppose every remaining constraint belongs to one block, and no permitted subsequent refinement introduces cross-block constraints. Then a global completion exists exactly when every block has a completion. If total completion cost is a sum of block costs, a minimum-cost global completion is assembled from minimum-cost block completions.
\end{proposition}
\begin{proof}
Restrict any global completion to each $V_i$ to obtain the forward implication. Conversely, take a completion for each block and unite the substitutions. Their domains are disjoint, their typing dependencies are local or fixed, and every constraint is satisfied within one block, so the union is well typed and satisfies all constraints. Additivity gives the minimum-cost assertion by replacing any nonminimal block of a proposed optimum.
\end{proof}

The final hypothesis about future refinement matters. A rule that later introduces a helper shared by two blocks changes their independence. Either forbid this after a block has been sealed, or rejoin the blocks. Similarly, a requirement to use a provider at most once, a global syntax-size budget, or a nonlinear ranking objective introduces additional coupling. A size budget can be handled by storing each component's cost-indexed solution stream and combining streams, rather than taking just one solution per component.

For cross-branch memoization, store a \emph{typed recipe}: an abstracted local telescope, normalized target, residual constraints, grammar identity, and a candidate term expressed over fresh formal parameters. Instantiate that recipe with fresh metavariables and recheck it in the receiving context. Do not cache a raw assignment containing an old \code{MVarId}. Negative cache entries additionally need the precise exhausted bound and a proof that the reused search space is covered. A failed bounded proof attempt is not a reusable impossibility fact.

This can be implemented incrementally. Keep today's whole-continuation rollback discipline. First use the graph only to prioritize informative goals. Next memoize closed successful recipes. Only then allow independently certified components to commit without global backtracking. This sequence earns performance without simultaneously changing every correctness argument.

\Q{R1.Q2}{An admissible remaining-cost bound}{What lower bound is correct when solving one goal can instantiate or close others?}

\answer The universally safe starting point is $h(S)=0$, with nonnegative fixed transition costs. Counting goals is not generally admissible. Consider two residual constraints $n=0$ and $n+1=1$ sharing an unknown $n$. In a cost model where assigning $n:=0$ costs one and reflexive propagation is free, one charged action closes both. A proposed bound of two exceeds the actual remaining cost of one. Giving exact-local steps zero cost creates additional counterexamples unless those steps are completely normalized before estimating cost.

Let $c^*(S)$ denote the least additional cost of any successful completion. For a lower bound $\ell_i$ attached to obligation $i$, the desired statement is $\ell_i\leq c^*(S)$, not merely that $\ell_i$ estimates the effort of solving that obligation in isolation. Within one coupled component,
\[
 h(S)=\max_i\ell_i
\]
is valid whenever each individual inequality is proved. Summing these bounds is not valid without a disjoint-charge argument. For components satisfying \cref{prop:factor}, with additive disjoint costs, independently established component bounds may be summed.

A more informative systematic approach is a \emph{relaxed search problem}. Erase selected constraints, add possible transitions, and assign each abstract transition no greater cost than the concrete transition it represents. If every concrete successful path projects to an abstract successful path, the abstract shortest-path distance is a lower bound. The abstraction must retain one-to-many propagation: charging separately for two goals that one concrete action discharges defeats admissibility. This method is a design framework, not a claim that a cheap, strong abstraction has already been found for Leant2.

With $h=0$, uniform-cost search is optimal only under the usual operational assumptions: the graph faithfully represents the intended completion space, all costs are fixed and nonnegative, and enough finite-cost work can be completed before the solution. Zero-cost cycles or infinitely many zero-cost alternatives can prevent progress. Use canonical administrative normalization plus deduplication, or a positive syntactic grade for every genuine expansion. If an admissible but inconsistent heuristic is later added to a graph search, allow reopening states when a better $g$ value is found.

\Q{R1.Q3}{Learning search costs without changing the promise}{Can just-in-time learning safely guide search and ranking, and how should budget effects be controlled?}

\answer Separate three quantities: a fixed program objective, accumulated charged work, and a learned proposal priority. They answer different questions. Probe demonstrates learning a synthesis distribution from partial successes; it does not establish that a deadline-limited learned run returns the same candidates as a fixed-cost exhaustive run.\cite{probe}

For Leant2, keep a stable objective such as a description length over the declared grammar. Use learned priorities to order a second frontier, or to break ties inside completed fixed-cost bands. Do not retrospectively decrease an edge cost while continuing to claim that the original A* proof applies. If the objective really changes, start a new optimization epoch, recompute affected bounds, and label the result with the new objective.

A deterministic portfolio can allocate a fixed fraction of logical work to a reference enumerator and the rest to adaptive search. Its important property is fairness, not a universal speedup theorem. If every finite reference prefix eventually receives work, no finite reference derivation is permanently starved. At a finite deadline, however, different priorities can still produce different answers. The receipt should therefore include the objective version, learning-state digest, ordered decisions, and completed reference bound.

Selection and final presentation should also remain separate. A learned heuristic can search aggressively for useful programs while final ranking applies a fixed, transparent ordering to the verified set. That gives stable ranking \emph{of a given set}, not stability of the set itself. A best-candidate guarantee requires a frontier lower-bound certificate against the same final objective; a 400\,ms grace interval does not provide it.

\Q{R1.Q4}{Machine-independent answers under a real deadline}{Can wall-clock-limited results be made independent of machine speed without a calibrated node budget?}

\answer Not in the unrestricted form suggested. A deterministic logical trace has prefixes $S_0,S_1,\ldots$. Place the first new candidate after step $j$. Two machines can complete respectively $j-1$ and $j$ steps before the same deadline. Unless the algorithm withholds the additional candidate or disregards the deadline, their results differ. Arbitrarily slow machines make any nontrivial fixed-deadline guarantee impossible.

The attainable guarantee is \emph{prefix determinism}. Specify a deterministic logical work sequence, return the exact completed prefix, and allow continuation from it. Charge at stable points such as rule applications, unification attempts, and verified observation reductions. A single ``node'' is often too coarse: a node can invoke a very expensive normalizer or prover. Maintain separate logical counters and physical time, with interruptible subbudgets for costly services. Wall time decides where the trace stops, not what that trace means.

For parallel execution, give tasks stable identifiers and deterministic work allotments. Commit completed epochs in task order. At a deadline, report the last fully committed epoch and explicitly separate opportunistic extra results. Sorting all candidates that happened to finish cannot undo the race that determined set membership. Similarly, resuming three two-second runs need not reproduce one six-second run, because replay and process startup consume time. Resuming to the same \emph{logical work prefix} can reproduce it.

Luby-style restart schedules have a theorem for repeated runs of Las Vegas algorithms under a runtime-distribution model. They do not automatically optimize heterogeneous deterministic synthesis lanes; restarting the same deterministic prefix repeatedly can simply waste work.\cite{luby} Use a resumable, fair geometric allocation as the default, and treat restart randomization as a separately measured policy. Recent snapshotting work concerns avoiding expensive re-elaboration in external tactic-search pipelines; its cost regime is not the same as Leant2's already-persistent in-process state.\cite{snapshot,evidence}
